\pdfinfoomitdate=1
\pdftrailerid{}
\pdfsuppressptexinfo=15
\newcommand{\DoNotLoadEpstopdf}{}
\documentclass[11pt]{article}
\usepackage[T1]{fontenc}
\usepackage[utf8]{inputenc}
\usepackage{lmodern,amsmath,amssymb,amsthm,mathtools,booktabs,array}
\usepackage[margin=1in]{geometry}
\usepackage{microtype}
\usepackage[hidelinks]{hyperref}
\usepackage{enumitem}
\setlist{itemsep=3pt,topsep=5pt}
\allowdisplaybreaks[2]
\emergencystretch=2em
\newtheorem{theorem}{Theorem}[section]
\newtheorem{lemma}[theorem]{Lemma}
\newtheorem{proposition}[theorem]{Proposition}
\newtheorem{corollary}[theorem]{Corollary}
\theoremstyle{remark}\newtheorem{remark}[theorem]{Remark}
\newcommand{\eps}{\varepsilon}
\newcommand{\OO}{\mathcal O}
\newcommand{\R}{\mathbb R}
\newcommand{\coef}[2]{[#1^{#2}]}
\DeclareMathOperator{\Li}{Li}
% Added by the ProveIt write phase (batch 110, 7 October 2026): dated notes
% marked [write], a macro for file, label and declaration names, longtable.
\usepackage{longtable}
\newcommand{\file}[1]{\nolinkurl{#1}}
\expandafter\def\expandafter\UrlBreaks\expandafter{\UrlBreaks\do\-\do\_\do\.}
\newenvironment{writenote}[1][]{\par\smallskip\begingroup\small
 \noindent\textbf{[write]}\if\relax\detokenize{#1}\relax\else~\textbf{#1.}\fi\ }%
 {\par\endgroup\smallskip}
\hypersetup{pdftitle={Report186 Rounded running means},pdfauthor={Research report},pdfsubject={OEIS A065094 and A065095 amplitude selection all orders asymptotics and inverses}}
\title{Report186\\[5pt]Rounded running means\\[3pt]OEIS A065094 and A065095}
\author{Exact amplitude selection and all order asymptotics}
\date{3 October 2026}
\begin{document}
\maketitle
\begin{abstract}
The floor and ceiling running-mean sequences have the modified-Bessel growth conjectured in their OEIS entries. More generally, every bounded additive perturbation of the unrounded recurrence has the form
\[
 a_n=A L_{n-1}(-1)+\OO(\sqrt n).
\]
We give an exact positive-weight formula for $A$, an exact pointwise optimal error interval for the full bounded-perturbation class, and nested rational bounds for $A$ from any finite trajectory. The latter certify the published numerical floor and ceiling constants without assumptions about fractional parts. Classical Perron asymptotics then transfer to every fixed order, with an exponentially smaller relative perturbation error. We give a rational coefficient generator, prove eventual strict log-concavity, and derive all fixed orders of the inverse along the exact sequence range. Arbitrary real thresholds require two monotone comparison functions and two ceilings; a single smooth approximation cannot remove the staircase ambiguity. The accompanying source archive includes exact standard-library certificates and finite symbolic algorithms, independent positive-sum replay, corruption checks, and a deterministic offline PDF builder.
\end{abstract}

\paragraph{Scope and attribution.}
The unrounded Laguerre solution, already recorded in A376995, and the classical Perron expansion are established inputs. The numerical constants in the target entries are credited to V\'aclav Kot\v e\v sovec in October 2024. The contribution developed here is their rigorous selection and certification for the rounded recurrences, together with sharp bounded-forcing control and its consequences. No worldwide novelty determination, proof-assistant certification, effective asymptotic onset, or distributional assumption about rounding errors is claimed.

\begin{writenote}[Status and trust boundary, added 7 October 2026, batch~110]
This is bundle Report~186 of an external AI-assisted research session, filed
in ProveIt as the single-source research report
\file{a065094-rounded-running-means} (provenance, the sources as the write
read them, what was checked, relation to the repository, the collected
non-claims and the reading conventions are in
Section~\ref{rrm:sec:provenance}). It is unrefereed and not formalized. Its
author line reads ``Exact amplitude selection and all order asymptotics'' and
its PDF author field ``Research report'': it names no person or tool.
\emph{What rests on what.} Theorems~\ref{rrm:thm:shadow}
and~\ref{rrm:thm:certificate} are proved by hand from exact identities, with
the classical Perron expansion of $L_m(-1)$ (credited to
Dea\~no--Huertas--Marcell\'an's statement) as the one special-function input;
Theorem~\ref{rrm:thm:allorders}, Proposition~\ref{rrm:prop:logconcavity} and
Theorems~\ref{rrm:thm:rangeinverse}, \ref{rrm:thm:threshold} follow by
conventional arguments. The 70-place amplitude intervals of
Section~\ref{rrm:sec:numerics} are exact rational computations of the
shipped standard-library programs. \emph{Two OEIS conjectures proved.}
Corollary~\ref{rrm:cor:oeis} proves the modified-Bessel growth asked for in
A065094 and A065095 (Remark~\ref{rrm:rem:oeis}), and every decimal of
Kot\v{e}\v{s}ovec's numerical constants there is confirmed as a truncation of
the certified values. The write also adds Remark~\ref{rrm:rem:transseries}
(the inverses against the transseries volume) and dated notes. No statement,
proof or number of the source was changed.
\end{writenote}

\newpage
{\small\tableofcontents}
\newpage

\section{Definitions and principal results}\label{rrm:sec:main}
Let $a_n^-$ and $a_n^+$ denote the sequences in OEIS A065094 and A065095, respectively \cite{oeisfloor,oeisceil}. With $S_n^\pm=\sum_{j=1}^n a_j^\pm$, their definitions are
\begin{align}
 a_1^-&=a_1^+=1,\notag\\
 a_{n+1}^-&=a_n^-+\left\lfloor\frac{S_n^-}{n}\right\rfloor,
 &a_{n+1}^+&=a_n^++\left\lceil\frac{S_n^+}{n}\right\rceil.
 \label{rrm:eq:rounded}
\end{align}
The initial values are
\[
\begin{array}{c|rrrrrrrrrr}
n&1&2&3&4&5&6&7&8&9&10\\\hline
 a_n^-&1&2&3&5&7&10&14&20&27&36\\
 a_n^+&1&2&4&7&11&16&23&33&46&62
\end{array}
\]
All terms are positive integers, and both sequences are strictly increasing: every running mean is at least one, so every rounded increment is at least one.

It is useful to treat a broader linear recurrence with arbitrary bounded forcing:
\begin{equation}\label{rrm:eq:forced}
 a_1=\alpha,\qquad S_n=\sum_{j=1}^n a_j,\qquad
 a_{n+1}=a_n+\frac{S_n}{n}+\eps_n,\qquad \ell\le\eps_n\le u.
\end{equation}
Here $\alpha,\ell,u\in\R$, $\ell\le u$, and no independence, periodicity, or limiting distribution is required. For floor rounding $\eps_n=-\{S_n/n\}\in(-1,0]$; for ceiling rounding $\eps_n\in[0,1)$. We may use the closed intervals $[-1,0]$ and $[0,1]$ in all bounds.

Define the positive sequences
\begin{align}
 P_n&=L_{n-1}(-1)=\sum_{j=0}^{n-1}\frac{\binom{n-1}{j}}{j!},
 &T_n&=\sum_{j=1}^n P_j,\label{rrm:eq:P}\\
 Q_n&=\int_0^\infty e^{-s}\frac{s^{n-1}}{(1+s)^n}\,ds,
 &w_n&=n(Q_n-Q_{n+1}).\label{rrm:eq:Q}
\end{align}
Set $H_1=0$ and $H_n=(n-1)P_{n-1}Q_n$ for $n\ge2$.

\begin{theorem}[Exact amplitude and shadowing]\label{rrm:thm:shadow}
For every trajectory \eqref{rrm:eq:forced}, the weights satisfy $w_n>0$ and $\sum_{n\ge1}w_n=1$. The absolutely convergent amplitude
\begin{equation}\label{rrm:eq:A}
 A=\alpha+\sum_{k\ge1}w_k\eps_k
\end{equation}
satisfies the exact identity
\begin{equation}\label{rrm:eq:shadow}
 a_n-AP_n
 =Q_n\sum_{k<n}k\eps_k(P_{k+1}-P_k)
  -P_n\sum_{k\ge n}w_k\eps_k.
\end{equation}
Consequently, for every $n\ge1$,
\begin{equation}\label{rrm:eq:sharp}
 \ell H_n-u(H_n+1)\le a_n-AP_n
 \le uH_n-\ell(H_n+1).
\end{equation}
Each endpoint is attainable at a prescribed $n$ by some forcing sequence in $[\ell,u]$. Moreover $H_n\sim\sqrt n/2$, and therefore
\begin{equation}\label{rrm:eq:Osqrt}
 a_n=AP_n+\OO(\sqrt n),
\end{equation}
uniformly over all forcings in the fixed interval $[\ell,u]$.
\end{theorem}

The error bound does not depend on $\alpha$. Pointwise sharpness refers to arbitrary bounded forcing; the extremizing forcing can depend on the observation index. We do not claim that either specific rounded orbit attains an error of order $\sqrt n$.

\begin{theorem}[Nested rational amplitude bounds]\label{rrm:thm:certificate}
For every $N\ge1$,
\begin{equation}\label{rrm:eq:bracket}
 L_N:=\frac{S_N+N\ell}{T_N}\le A\le
 U_N:=\frac{S_N+Nu}{T_N}.
\end{equation}
The lower endpoints are nondecreasing, the upper endpoints are nonincreasing, and their exact width is
\begin{equation}\label{rrm:eq:width}
 U_N-L_N=\frac{N(u-\ell)}{T_N}
 \sim 2\sqrt\pi e^{1/2}(u-\ell)N^{3/4}e^{-2\sqrt N}
\end{equation}
when $u>\ell$. For $u=\ell$ the width is identically zero. Rational finite input gives a rational certificate.
\end{theorem}

\begin{corollary}[The two OEIS growth conjectures]\label{rrm:cor:oeis}
There are positive constants $A_-$ and $A_+$ such that, for either choice of sign,
\begin{align}
 a_n^\pm&=A_\pm P_n+\OO(\sqrt n),\label{rrm:eq:bothshadow}\\
 a_n^\pm&\sim C_\pm I_0(2\sqrt n)
 \sim c_\pm n^{-1/4}e^{2\sqrt n},\label{rrm:eq:bessel}
\end{align}
where
\begin{equation}\label{rrm:eq:normalization}
 C_\pm=A_\pm e^{-1/2},\qquad
 c_\pm=\frac{A_\pm e^{-1/2}}{2\sqrt\pi}.
\end{equation}
The published numerical constants are confirmed by exact rational enclosures in Section~\ref{rrm:sec:numerics}. Both sequences share every normalized algebraic correction from the classical Laguerre expansion.
\end{corollary}

The rest of the report proves these statements and makes the all-order and inverse conclusions precise. The only all-order special-function input is the classical Perron expansion, stated in Section~\ref{rrm:sec:perron}; all perturbation identities and amplitude bounds are proved directly.

\begin{remark}[{[write]} The OEIS entries, quoted, 7~October 2026]\label{rrm:rem:oeis}
The entries were read on 7~October 2026 in the OEIS internal format; the
quotations are verbatim.
\begin{enumerate}[label=(\alph*)]\raggedright
\item A065094~\cite{oeisfloor} (revision \#19, 12~October 2024; author Ulrich
 Schimke), ``\texttt{a(1) = 1, a(n+1) is the sum of a(n) and floor( arithmetic
 mean of a(1) ... a(n) ).}'', has the unsigned comment
 \begin{center}\small
 \parbox{0.92\textwidth}{\raggedright\texttt{It seems that a(n) is asymptotic to C*BesselI(0,2*sqrt(n)) where C is a constant C = 0.44... and BesselI(b,x) is the modified Bessel function of the first kind. Can someone prove this?}}
 \end{center}
 followed by Kot\v{e}\v{s}ovec's comment (12~October 2024) ``\texttt{Numerically,
 a(n) \textasciitilde{} c * exp(2*sqrt(n)) / n\^{}(1/4), where c =
 0.12571987512700920098166979884420897638511306007242...}'' and ``\texttt{C =
 0.445665353608456118285630970456186510059368576678...}''. Its comment
 ``\texttt{A241772(n) = a(n+1) - a(n) = (Sum\_\{1..n\} a(k)) / n.}''
 (Zumkeller) omits the floor; this is the transcription inconsistency named in
 Section~\ref{rrm:sec:sources}.
\item A065095~\cite{oeisceil} (revision \#20, 30~September 2026; Schimke),
 the ceiling analogue, has the same unsigned question with ``\texttt{C =
 0.78...}'' and Kot\v{e}\v{s}ovec's ``\texttt{c =
 0.2214496835182522607818590241239262909281832289078...}'', ``\texttt{C =
 0.78501868866746800511978860290796656518270697588...}''. Its example
 ``\texttt{a(5) = a(4) + ceiling((a(1)+a(2)+a(3)+a(4))/4) = 7 + ceiling((1+2+4+7)/4) = 7 + floor(14/4) = 7 + 4 = 11.}''
 writes \texttt{floor} where the value is the ceiling.
\item A376995~\cite{oeisunrounded} (revision \#16, Kot\v{e}\v{s}ovec) is
 $\lfloor b(n)\rfloor$ for the unrounded recurrence; it gives
 ``\texttt{a(n) \textasciitilde{} exp(-1/2) * BesselI(0, 2*sqrt(n))}'' and
 ``\texttt{a(n) = floor(A160617(n-1)/A160618(n-1))}'', where A160617 and
 A160618 are the numerator and denominator of $L_n(-1)$: so $b(n)=P_n$.
\end{enumerate}
\emph{What is proved.} Corollary~\ref{rrm:cor:oeis} proves the two unsigned
statements exactly as asked: $a_n^\pm\sim C_\pm I_0(2\sqrt n)$ with
$C_-=0.4456\ldots$ and $C_+=0.7850\ldots$, and Kot\v{e}\v{s}ovec's numerical
equivalents $a_n^\pm\sim c_\pm e^{2\sqrt n}n^{-1/4}$. His four decimals (50,
48, 49 and 47 places) are truncations of the certified values of
Section~\ref{rrm:sec:numerics}, as are the source's 48-place
values~\eqref{rrm:eq:decimalminus}--\eqref{rrm:eq:decimalplus}. Both b-files,
fetched on 7~October 2026, are byte-identical to the shipped
\file{data/b065094.txt} and \file{data/b065095.txt}, and the write's own
integer recomputation of both orbits agrees with all $1000$ terms of each.
OEIS content is published under CC~BY-SA~4.0. Nothing was submitted to the
OEIS.
\end{remark}

\subsection{Provenance, sources and reading conventions}\label{rrm:sec:provenance}
\begin{writenote}[Added 7 October 2026, batch~110]
\emph{Provenance.} This report is Report~186 of the session bundle that
ProveIt's \file{docs/incoming} intake processed as batch~110: the archive
\file{Rounded_Running_Means_Bessel_Asymptotics_and_Inverses_Source.zip}
($617{,}579$ bytes, SHA-256 \file{5a8332939ea7...8cea55487d0ef1},
33 files, no wrapper directory), arrival commit \file{60f54ea06}, placed in
\file{8622ca7e5}. The text is the delivered \file{Report186.tex} (760 lines,
dated 3~October 2026; the delivered PDF has 19 pages), shipped as
\file{article.tex}. The programs are shipped under \file{code/} (the
delivered root \file{build.py}, \file{reproduce.py}, \file{test_build.py}
and \file{verify_manifest.py} among them; the optional ones as
\file{code/optional-*.py}), the certificates as
\file{data/certificates-*.json}, the generated summaries as
\file{data/generated-*.json}, the two OEIS b-files, the prefixes and the
source hashes under \file{data/}, and \file{DATA_SOURCES.md} and
\file{README_REPRODUCIBILITY.md} at the report root (the report README maps
every name). Not shipped, and retrievable from the arrival commit: the
delivered PDF, the checksum manifest \file{SHA256SUMS.json} (32 entries,
verified at placement) and the delivered README, replaced by the report
README. The package records no ProveIt commit and names no repository path;
it carries no ``prepared for private review'' line and no personal data. It
is a single-source report, so no merge choice was made. The write prefixed
the 101 delivered labels with \file{rrm:} and updated the 63 references to
them; it added this subsection, the notes marked [write],
Remarks~\ref{rrm:rem:oeis} and~\ref{rrm:rem:transseries}, and dated notes.
The added remarks are the last statements of their sections, this
subsection follows the last delivered text of Section~\ref{rrm:sec:main},
and the added displays are unnumbered, so every theorem, section and
equation of the source keeps its delivered number.

\emph{The sources, as the write read them} (7~October 2026).
\begin{itemize}\raggedright
\item OEIS A065094, A065095 with their b-files, A376995, and the names of
 A160617 and A160618, in the internal format (Remark~\ref{rrm:rem:oeis}).
\item The transseries volume
 \path{Analysis/Transseries/docs/series-and-transseries/Transseries_And_Inversion/transseries_and_inversion.tex}:
 \file{p0:thm:lambert-core}, \file{p0:thm:core-reversion},
 \file{p0:def:three-inverses} and \file{p0:thm:staircase}, read for
 Remark~\ref{rrm:rem:transseries}.
\item Not read by the write: Dea\~no, Huertas and Marcell\'an~\cite{dhm}
 (the source's reading of its formula~(3) and formulas (24)--(27) is
 reported, not confirmed; the write recomputed the two coefficients the
 source takes from it, see below) and the DLMF~\cite{dlmf}.
\end{itemize}

\emph{What was checked at intake.} The batch-110 dossier read the
manuscript and found no error. On copies, \file{reproduce.py --output-dir}
passed in 13~s; the normal-mode certificates were byte-identical to the six
recorded \file{certificates/*.json}, and its result equal to
\file{generated/verification.json} modulo CR. \emph{What the write checked}
(Python~3.14.4, SymPy~1.14.0, own code). Every proof was read line by line,
with no error found; in particular the impulse response~\eqref{rrm:eq:G}
(both initial values), the tail~\eqref{rrm:eq:tail}, the
envelopes~\eqref{rrm:eq:floorenv}--\eqref{rrm:eq:ceilenv}, the positivity
bound~\eqref{rrm:eq:positivity}, the width
equivalent~\eqref{rrm:eq:width}, the ratio expansion~\eqref{rrm:eq:ratios},
and the constant-forcing example of Section~\ref{rrm:sec:questions}. With its
own SymPy code it rederived $c_1,\dots,c_6$ of~\eqref{rrm:eq:ctable} from
the recurrence~\eqref{rrm:eq:recP} by the ansatz of~\eqref{rrm:eq:Pall}, the
$b_j$ of~\eqref{rrm:eq:bvalues}, the five coefficients
of~\eqref{rrm:eq:delta} from~\eqref{rrm:eq:deltagen}, and the
range inverse~\eqref{rrm:eq:rangeinverse}; and it checked
$c_1=d_1-1$, $c_2=d_2-d_1+3/4$ from the quoted $d_1=31/48$, $d_2=169/4608$:
all agree. With exact integers it recomputed both orbits and, at $N=3000$,
the brackets~\eqref{rrm:eq:floorcert}--\eqref{rrm:eq:ceilcert} with $T_N$ from
the positive sum~\eqref{rrm:eq:positiveT}: they agree with the table of
Section~\ref{rrm:sec:numerics} in all $40$ places printed by the check. The
note at the end of Section~\ref{rrm:sec:questions} records a residual
diagnostic. Floating computations are observations, not certificates.

\emph{Relation to the repository.} Before batch~110 no file of the
repository named A065094, A065095 or A376995. No statement of this report is
formalized in Lean or Rocq, and its place in the collection confers no
formal status. No neighbouring report treats rounded recurrences or Laguerre
amplitudes, so no reciprocal note is proposed. The inverses are compared with
the transseries volume in Remark~\ref{rrm:rem:transseries}.
\end{writenote}

\begin{writenote}[What is not claimed, collected from the source]
The unrounded Laguerre solution (A376995), the classical Perron expansion and
its coefficient methods, and the numerical constants (Kot\v{e}\v{s}ovec, 2024)
are prior; the source claims ``No worldwide novelty determination,
proof-assistant certification, effective asymptotic onset, or distributional
assumption about rounding errors'', and its source review ``is not an
exhaustive literature search and does not establish global priority''. The
optimal envelope is sharp only for the whole bounded-forcing class: it is
``not a lower bound on either actual rounded residual''. ``All orders'' means
one theorem for each fixed order, with no convergence of the formal series
and no uniformity in the order. Agreement with the published constants is a
consistency check; the proof is the exact inequalities. Finite and floating
observations are not onset certificates; no effective onset is given for the
all-order remainder, log-concavity, nearest-integer recovery or the threshold
comparison functions. The range inverse needs the exact amplitude; the
threshold theorem is asymptotic, not an executable finite certificate, and a
single ceiling of a truncation is not asserted to give the threshold. The
write adds: its checks are exact finite or floating computations; it claims
no novelty for any inversion.
\end{writenote}

\medskip
\begingroup\small
\noindent\emph{Reading conventions} (letters the source reuses, with the
tempting false readings; no symbol of the source was renamed; symbols of the
transseries volume carry the subscript ``vol'' in
Remark~\ref{rrm:rem:transseries}).\par\smallskip
\renewcommand{\arraystretch}{1.1}
\begin{longtable}{@{}>{\raggedright\arraybackslash}p{0.95in}>{\raggedright\arraybackslash}p{3.3in}>{\raggedright\arraybackslash}p{1.85in}@{}}
\toprule
Symbol & Meaning here (where) & Not to be read as\\
\midrule
\endhead
$A$, $A_\pm$ & the amplitude~\eqref{rrm:eq:A} and its floor and ceiling values & the volume's $A_n$, $\mathcal A_{\mathrm{vol}}$\\
$C_\pm$, $c_\pm$, $c_j$, $c$ & the Bessel and exponential constants~\eqref{rrm:eq:normalization}; the coefficients~\eqref{rrm:eq:ctable}; $c$ for $c_\pm$ in Sections~\ref{rrm:sec:inverse}--\ref{rrm:sec:threshold}; $C(h)$ the series of Section~\ref{rrm:sec:coefficients} & each other\\
$L_m$, $L_N$, $L_M$ & Laguerre polynomials; the lower bracket~\eqref{rrm:eq:bracket}; the lower comparison function of Section~\ref{rrm:sec:threshold} & each other; a logarithm; the volume's $L_{\mathrm{vol}}$\\
$U_N$, $U_M$ & the upper bracket; the upper comparison function & each other\\
$P_n$, $Q_n$, $T_n$, $S_n$ & $L_{n-1}(-1)$; the minimal solution~\eqref{rrm:eq:Q}; $\sum_{j\le n}P_j$; partial sums of the orbit & each other\\
$H_n$, $h$ & $(n-1)P_{n-1}Q_n$; $n^{-1/2}$ & each other\\
$D$, $d_j$ & $c\sqrt2$~\eqref{rrm:eq:Dt}; the coefficients of the degree-$m$ Perron expansion~\eqref{rrm:eq:perrondegree} & each other; the volume's $d_{\mathrm{vol}}$\\
$\delta$, $d_r$ & the displacement $t-t_0$; its coefficients & the volume's $e_n$\\
$t$, $t_0$, $v$ & $2\sqrt n$; the Lambert root~\eqref{rrm:eq:tzero}; $1/t_0$ & the volume's $t$ (its $t$ plays the role of $v$ here)\\
$E_1$, $E_m$, $R_m$ & the exponential integral; partial sums of $e^{1/2}$ and their remainders (Section~\ref{rrm:sec:numerics}) & each other\\
$F_n$, $F_M$, $f_n$, $f_J$ & $(n-1)!P_n$; the carrier~\eqref{rrm:eq:carrier}; the integrand of $Q_n$; the truncated logarithmic forward function & each other\\
$G_{n,k}$, $W_n$, $w_n$ & the impulse response; the Casoratian; the weights & each other; Lambert's $W_{-1}$\\
$\beta_N$, $b_j$, $B_n$, $b_n$ & the determinant-like quantity~\eqref{rrm:eq:beta}; the logarithmic coefficients~\eqref{rrm:eq:bvalues}; shifted sums and terms in Section~\ref{rrm:sec:certproof} & each other; the volume's $b_{\mathrm{vol}}$\\
$\kappa$, $\ell$, $u$, $r$ & $e^{-1/2}/(2\sqrt\pi)$; the forcing bounds; a constant shift & the volume's $\kappa_{\mathrm{vol}}$\\
$N$, $N(Y)$, $M$, $K$ & a certificate index; the threshold~\eqref{rrm:eq:threshold}; a fixed order; the forward constant & each other\\
\bottomrule
\end{longtable}
\endgroup

\section{The growing and decaying solutions}\label{rrm:sec:fundamental}
\subsection{The Laguerre solution}
The classical Laguerre recurrence gives $P_1=1$, $P_2=2$, and
\begin{equation}\label{rrm:eq:recP}
 nP_{n+1}-2nP_n+(n-1)P_{n-1}=0\qquad(n\ge2).
\end{equation}
For convenience at $n=1$ set $P_0=0$; then \eqref{rrm:eq:recP} also holds there. Writing $\Delta P_n=P_{n+1}-P_n$ gives
\[
 n\Delta P_n-(n-1)\Delta P_{n-1}=P_n.
\]
Summing, with the initial value $P_2-P_1=1$, proves
\begin{equation}\label{rrm:eq:T}
 T_n=n(P_{n+1}-P_n).
\end{equation}
Thus $P_{n+1}=P_n+T_n/n$: $P$ is exactly the unrounded trajectory with initial value one. This identification is already part of the record A376995 \cite{oeisunrounded}. The finite sum in \eqref{rrm:eq:P} proves positivity directly. It also supplies an independent finite computation, separate from the recurrence.

\subsection{The positive minimal solution}
The integral \eqref{rrm:eq:Q} converges absolutely. Its integrand is positive, and multiplication by $s/(1+s)<1$ shows $0<Q_{n+1}<Q_n$. For $n\ge2$, differentiate
\[
 f_n(s)=e^{-s}\frac{s^{n-1}}{(1+s)^n}.
\]
The exact derivative is
\begin{equation}\label{rrm:eq:derivative}
 f_n'(s)=e^{-s}\frac{s^{n-2}}{(1+s)^{n+1}}
 \{n-1-2s-s^2\}.
\end{equation}
Both endpoint values of $f_n$ vanish. Combining the three integrands in the left side below produces the same numerator in braces, so integration gives
\begin{equation}\label{rrm:eq:recQ}
 nQ_{n+1}-2nQ_n+(n-1)Q_{n-1}=0\qquad(n\ge2).
\end{equation}
The special initial values are
\begin{equation}\label{rrm:eq:Qinitial}
 Q_1=eE_1(1),\qquad Q_2=2Q_1-1,
 \qquad E_1(1)=\int_1^\infty e^{-t}\frac{dt}{t}.
\end{equation}
For the second identity, integrate the derivative of $e^{-s}/(1+s)$ and use
$s/(1+s)^2=1/(1+s)-1/(1+s)^2$. Also $Q_1<\int_0^\infty e^{-s}ds=1$.

\begin{lemma}[Normalized Casoratian]\label{rrm:lem:casoratian}
For every $n\ge1$,
\begin{equation}\label{rrm:eq:casoratian}
 n(P_{n+1}Q_n-P_nQ_{n+1})=1.
\end{equation}
\end{lemma}
\begin{proof}
Put $W_n=P_{n+1}Q_n-P_nQ_{n+1}$. Substitution of the two recurrences gives $nW_n=(n-1)W_{n-1}$ for $n\ge2$. At $n=1$, \eqref{rrm:eq:Qinitial} gives $W_1=2Q_1-Q_2=1$. This proves the identity and linear independence.
\end{proof}

The adjective minimal will be justified by $Q_n/P_n\to0$ and a full leading equivalent below. It does not require a numerical backward recurrence, which would introduce avoidable stability questions into the certificate.

\section{Positive weights and exact Green formula}\label{rrm:sec:green}
\subsection{Mass and tails}
Subtract the integrands for $Q_n$ and $Q_{n+1}$ to obtain
\begin{equation}\label{rrm:eq:wint}
 w_n=n\int_0^\infty e^{-s}\frac{s^{n-1}}{(1+s)^{n+1}}\,ds>0.
\end{equation}
For $r=s/(1+s)$, the geometric derivative identity
$\sum_{n\ge1}nr^{n-1}=(1-r)^{-2}$ and Tonelli's theorem give
\begin{equation}\label{rrm:eq:mass}
 \sum_{n\ge1}w_n=\int_0^\infty e^{-s}\,ds=1.
\end{equation}
An integration by parts in $e^{-s}(s/(1+s))^n$ also gives
\begin{equation}\label{rrm:eq:walt}
 w_n=\int_0^\infty e^{-s}\left(\frac{s}{1+s}\right)^n ds.
\end{equation}
Hence the series defining $A$ is absolutely convergent, and
\begin{equation}\label{rrm:eq:roughA}
 \alpha+\ell\le A\le\alpha+u.
\end{equation}

For $n\ge2$ two exact sums will be needed:
\begin{align}
 \sum_{k<n}k(P_{k+1}-P_k)&=(n-1)P_{n-1},\label{rrm:eq:past}\\
 \sum_{k\ge n}w_k&=(n-1)Q_{n-1}.\label{rrm:eq:tail}
\end{align}
To prove \eqref{rrm:eq:past}, summation by parts gives $(n-1)P_n-T_{n-1}$, which is the claimed expression by \eqref{rrm:eq:T}. For \eqref{rrm:eq:tail}, Tonelli applied to \eqref{rrm:eq:walt} gives
\[
 \sum_{k\ge n}w_k
 =\int_0^\infty e^{-s}\frac{s^n}{(1+s)^{n-1}}\,ds.
\]
This equals $(n-1)Q_{n-1}$. Indeed the integral of the derivative of
$e^{-s}s^{n-1}/(1+s)^{n-2}$ is zero; its numerator over $(1+s)^{n-1}$ is
$(n-1)s^{n-2}-s^n$. The endpoint terms vanish for $n\ge2$. This argument establishes the tail without any unproved limiting cancellation.

Equations \eqref{rrm:eq:past}, \eqref{rrm:eq:tail}, and the Casoratian at $n-1$ give
\begin{equation}\label{rrm:eq:massdiff}
 P_n\sum_{k\ge n}w_k-Q_n\sum_{k<n}k(P_{k+1}-P_k)=1.
\end{equation}
At $n=1$ it follows from $P_1=1$ and \eqref{rrm:eq:mass}. The positive past coefficient total is $H_n$ and the positive future coefficient total is $H_n+1$.

\subsection{An impulse with the correct initial jump}
Eliminating $S_n$ from \eqref{rrm:eq:forced} yields, with $\eps_0=0$,
\begin{equation}\label{rrm:eq:secondorder}
 na_{n+1}-2na_n+(n-1)a_{n-1}
 =n\eps_n-(n-1)\eps_{n-1}.
\end{equation}
At $n=1$ the term involving $a_0$ vanishes. A unit error at time $k$ has zero response for $n\le k$. For $n>k$ define
\begin{equation}\label{rrm:eq:G}
 G_{n,k}=w_kP_n+k(P_{k+1}-P_k)Q_n.
\end{equation}
The same expression, if evaluated at $n=k$ and $n=k+1$, equals one in both cases by \eqref{rrm:eq:casoratian}. Its homogeneous extension consequently has
\[
 G_{k+2,k}=\frac{k+2}{k+1}.
\]
These are exactly the two actual initial response values: the error adds one to $a_{k+1}$ and to $S_{k+1}$, then the next running mean adds $1/(k+1)$. Thereafter the second-order recurrence is homogeneous. Thus \eqref{rrm:eq:G}, together with zero response through time $k$, is the correct impulse solution. This check is important: differencing the forcing creates the additional negative impulse at the next step in \eqref{rrm:eq:secondorder}.

By finite superposition,
\begin{equation}\label{rrm:eq:variation}
 a_n=\alpha P_n+\sum_{k<n}\eps_kG_{n,k}.
\end{equation}
Subtracting \eqref{rrm:eq:A} times $P_n$ proves \eqref{rrm:eq:shadow}. Only finite sums and the absolutely convergent amplitude tail are used.

\subsection{Sharpness of the envelope}
Every coefficient of a past error in \eqref{rrm:eq:shadow} is positive, and every coefficient of a future error is negative. Their total magnitudes are $H_n$ and $H_n+1$. This proves \eqref{rrm:eq:sharp}. For the upper endpoint choose $\eps_k=u$ for $k<n$ and $\eps_k=\ell$ for $k\ge n$; reverse the choices for the lower endpoint. These are legitimate bounded forcings and prove exact pointwise optimality, including $n=1$.

A useful specialization is
\begin{align}
 -H_n\le a_n^--A_-P_n&\le H_n+1,\label{rrm:eq:floorenv}\\
 -(H_n+1)\le a_n^+-A_+P_n&\le H_n.\label{rrm:eq:ceilenv}
\end{align}
The common order $\sqrt n$ is a theorem about the available uniform error envelope, not a lower bound on either actual rounded residual.

\section{Classical asymptotics and bounded forcing transfer}\label{rrm:sec:perron}
\subsection{The classical input and the index shift}
We use the outer Perron expansion of Laguerre polynomials, as stated in Dea\~no, Huertas, and Marcell\'an \cite[formula (3)]{dhm}. At fixed $z=-1$ and parameter zero, for each fixed integer $M\ge0$ it gives
\begin{equation}\label{rrm:eq:perrondegree}
 L_m(-1)=\frac{e^{-1/2}}{2\sqrt\pi}\,m^{-1/4}e^{2\sqrt m}
 \left\{\sum_{j=0}^M d_jm^{-j/2}
       +\OO_M(m^{-(M+1)/2})\right\},\quad d_0=1.
\end{equation}
The negative real point $-1$ lies in its permitted domain off the nonnegative real axis. This is an asymptotic theorem with a remainder for every fixed order; a formal recurrence alone would not establish it. The paper credits the classical result to Perron and Szeg\H{o}, and gives established coefficient constructions and checks.

Our index is $n=m+1$. Expanding the elementary prefactor and the finite series in \eqref{rrm:eq:perrondegree} gives
\begin{equation}\label{rrm:eq:Pall}
 P_n=\kappa n^{-1/4}e^{2\sqrt n}
 \left\{\sum_{j=0}^M c_jn^{-j/2}
       +\OO_M(n^{-(M+1)/2})\right\},\qquad
 \kappa=\frac{e^{-1/2}}{2\sqrt\pi}.
\end{equation}
Thus $P_n\sim\kappa n^{-1/4}e^{2\sqrt n}$ and $P_{n-1}/P_n\to1$. The degree-to-index shift matters: formula (27) in \cite{dhm} gives $d_1=31/48$ and $d_2=169/4608$, whereas the coefficients for $P_n$ are $c_1=-17/48$ and $c_2=649/4608$.

\subsection{The minimal solution without a second special function theorem}
Divide \eqref{rrm:eq:casoratian} by $nP_nP_{n+1}$ and sum. Since $0<Q_n\le Q_1$ and $P_n\to\infty$, the result is the exact positive tail
\begin{equation}\label{rrm:eq:reductionorder}
 \frac{Q_n}{P_n}=\sum_{j\ge n}\frac1{jP_jP_{j+1}}.
\end{equation}
The leading equivalent in \eqref{rrm:eq:Pall} implies, uniformly in all $j\ge n$ as $n\to\infty$,
\[
 \frac1{jP_jP_{j+1}}
 \sim\kappa^{-2}j^{-1/2}e^{-4\sqrt j}.
\]
Uniformity here follows directly from the definition of a convergent relative error on a tail of a sequence; no growing-parameter version of Perron's theorem is needed.

The positive function $f(x)=x^{-1/2}e^{-4\sqrt x}$ is decreasing and
\[
 \int_n^\infty f(x)\,dx=\frac12 e^{-4\sqrt n},
 \qquad \frac{f(n)}{\int_n^\infty f(x)\,dx}=\frac2{\sqrt n}.
\]
Integral comparison and \eqref{rrm:eq:reductionorder} therefore give
\begin{equation}\label{rrm:eq:Qasymp}
 Q_n\sim\frac1{2\kappa}n^{-1/4}e^{-2\sqrt n}
 =\sqrt\pi e^{1/2}n^{-1/4}e^{-2\sqrt n}.
\end{equation}
In particular $nQ_n\to0$, $Q_n/P_n\to0$, and
\begin{equation}\label{rrm:eq:Hasymp}
 H_n=(n-1)P_{n-1}Q_n\sim\frac{\sqrt n}{2}.
\end{equation}
Together with the exact interval, this finishes the proof of Theorem~\ref{rrm:thm:shadow}.

\subsection{Positivity and the Bessel conjectures}
For floor rounding, mass one rewrites the amplitude as
\[
 A_-=\sum_{k\ge1}w_k(1+\eps_k^-)>0.
\]
Every summand is nonnegative and $\eps_1^-=0$, so in fact
\begin{equation}\label{rrm:eq:positivity}
 A_-\ge w_1=1-Q_1=1-eE_1(1)>0.
\end{equation}
For ceiling rounding, $A_+\ge1$. Thus both leading amplitudes are genuinely nonzero; the generic inequality $\alpha+\ell\le A$ alone would not have proved this for floor rounding.

For fixed $A\ne0$, the relative shadowing error is
\begin{equation}\label{rrm:eq:relative}
 \frac{a_n-AP_n}{AP_n}=\OO\left(n^{3/4}e^{-2\sqrt n}\right).
\end{equation}
The standard leading modified-Bessel expansion \cite[10.40.1]{dlmf} is
$I_0(x)\sim e^x/\sqrt{2\pi x}$ as $x\to+\infty$. Applying it at $x=2\sqrt n$ and comparing with \eqref{rrm:eq:Pall} proves Corollary~\ref{rrm:cor:oeis} with the constants in \eqref{rrm:eq:normalization}.

\begin{theorem}[Every fixed forward order]\label{rrm:thm:allorders}
For either rounded sequence and each fixed $M\ge0$,
\begin{equation}\label{rrm:eq:allorders}
 a_n^\pm=c_\pm e^{2\sqrt n}n^{-1/4}
 \left\{\sum_{j=0}^M c_jn^{-j/2}
       +\OO_M(n^{-(M+1)/2})\right\}.
\end{equation}
The rational coefficients $c_j$ are identical for floor and ceiling rounding and coincide with those of $L_{n-1}(-1)$ in \eqref{rrm:eq:Pall}. The bounded-forcing relative error is smaller than every algebraic power.
\end{theorem}
\begin{proof}
The ratio in \eqref{rrm:eq:relative} is $o(n^{-B})$ for every fixed $B>0$, so it can be absorbed into the remainder in \eqref{rrm:eq:Pall} for any fixed $M$.
\end{proof}

All orders always means one theorem for each fixed truncation order. It does not mean that the infinite formal series converges or that the constants are uniform when the order grows with $n$.

\section{A rational generator for every coefficient}\label{rrm:sec:coefficients}
Let $h=n^{-1/2}$ and let $C(h)=\sum_{j\ge0}c_jh^j$ be a formal power series, normalized by $c_0=1$. For any formal series $C$, define
\begin{align}
 \mathcal R[C](h)={}&
 e^{2(\sqrt{1+h^2}-1)/h}(1+h^2)^{-1/4}
 C\left(\frac{h}{\sqrt{1+h^2}}\right)-2C(h)\notag\\
 &+(1-h^2)e^{2(\sqrt{1-h^2}-1)/h}(1-h^2)^{-1/4}
 C\left(\frac{h}{\sqrt{1-h^2}}\right).
 \label{rrm:eq:operator}
\end{align}
This is precisely \eqref{rrm:eq:recP} divided by $ne^{2\sqrt n}n^{-1/4}$. The apparent divisions by $h$ are harmless because $\sqrt{1\pm h^2}-1$ has order two. Binomial and exponential series therefore define every coefficient using finitely many rational operations.

For $j\ge1$, expansion through the first nonzero order gives
\begin{equation}\label{rrm:eq:triangular}
 \mathcal R[h^j]=-jh^{j+3}+\OO(h^{j+4}),
\end{equation}
while $\mathcal R[1]$ begins at degree four. A direct way to see the $-j$ coefficient is to compare the two substitutions $h^j(1\pm h^2)^{-j/2}$: the order-$h^2$ contributions cancel between the two shifts, and their products with the opposite order-$h$ exponential terms total $-jh^{j+3}$. The common $j$-independent contribution at degree $j+3$ vanishes. Consequently
\begin{equation}\label{rrm:eq:cgen}
 c_j=\frac1j\coef{h}{j+3}
 \mathcal R\left[\sum_{r=0}^{j-1}c_rh^r\right],\qquad j\ge1.
\end{equation}
The asymptotic existence supplied by Perron's theorem and this triangular uniqueness show that the generated series is the actual coefficient series. Formal cancellation is used to identify coefficients, not to prove the analytic remainder theorem.

The first seven coefficients are
\begin{equation}\label{rrm:eq:ctable}
\begin{array}{c|r}
j&c_j\\\hline
0&1\\[2pt]
1&-17/48\\[2pt]
2&649/4608\\[2pt]
3&-56533/3317760\\[2pt]
4&7946069/637009920\\[2pt]
5&937900373/42807066624\\[2pt]
6&2719850722091/308210879692800
\end{array}
\end{equation}
As an external normalization check, the first two coefficients in the degree-$m$ formula \eqref{rrm:eq:perrondegree} are obtained from \cite[formula (27)]{dhm}. The shift $m=n-1$ changes the elementary prefactor by
\[
 \frac{(n-1)^{-1/4}e^{2\sqrt{n-1}}}{n^{-1/4}e^{2\sqrt n}}
 =1-n^{-1/2}+\frac34n^{-1}+\OO(n^{-3/2}).
\]
Thus $c_1=d_1-1=-17/48$ and $c_2=d_2-d_1+3/4=649/4608$, in agreement with the recurrence generator.

The supplied mandatory implementation uses truncated arrays of rational numbers for addition, multiplication, substitution, exponential, logarithm, and rational binomial powers. Every fixed requested order is a finite computation. A separate logarithmic-ratio check verifies the displayed forward coefficients, and direct formal composition checks the displayed inverse coefficients. Optional symbolic-system scripts offer another route but are not required to reproduce the exact results.

\section{Nested certificates from the finite trajectory}\label{rrm:sec:certproof}
We now prove Theorem~\ref{rrm:thm:certificate} without evaluating $Q_n$ numerically. Fix a real constant $r$ and write
\[
 b_n=a_n+r,\qquad B_n=S_n+nr.
\]
Then $b_{n+1}=b_n+B_n/n+(\eps_n-r)$. Define the determinant-like quantity
\begin{equation}\label{rrm:eq:beta}
 \beta_N(r)=T_N(a_N+r)-P_N(S_N+Nr).
\end{equation}
At $N=1$ it is zero. Subtracting the two first-order systems for $b$ and $P$ gives
\begin{equation}\label{rrm:eq:betaupdate}
 \beta_{N+1}(r)=\beta_N(r)+T_N(\eps_N-r).
\end{equation}
Hence, using \eqref{rrm:eq:T},
\begin{equation}\label{rrm:eq:betasum}
 \beta_N(r)=\sum_{k<N}k(\eps_k-r)(P_{k+1}-P_k).
\end{equation}
A second direct subtraction gives
\begin{equation}\label{rrm:eq:ratioupdate}
 \frac{B_{N+1}}{T_{N+1}}-\frac{B_N}{T_N}
 =\frac{\beta_N(r)+T_N(\eps_N-r)}{T_NT_{N+1}}.
\end{equation}
For $r=\ell$, all terms in the numerator are nonnegative; for $r=u$, they are nonpositive. This proves nestedness of the proposed endpoints.

Their difference is exactly $N(u-\ell)/T_N$. To find its size, use the first two orders of \eqref{rrm:eq:Pall} in \eqref{rrm:eq:T} to obtain
\begin{equation}\label{rrm:eq:Tasymp}
 \frac{P_{N+1}}{P_N}=1+N^{-1/2}+\OO(N^{-1}),
 \qquad T_N\sim\sqrt N\,P_N.
\end{equation}
Therefore the width tends to zero and has the equivalent \eqref{rrm:eq:width} when $u>\ell$.

Both endpoints consequently tend to the same limit. To identify it, sum the exact shadowing estimate: $S_N=AT_N+\OO(N^{3/2})$. Since $T_N$ grows faster than every power, $S_N/T_N\to A$, and $Nr/T_N\to0$ for fixed $r$. Thus their common limit is exactly \eqref{rrm:eq:A}. Monotonicity then proves that every finite interval contains $A$, completing the theorem.

\subsection{The actual floor and ceiling brackets}
For the two trajectories the bounds become
\begin{align}
 \frac{S_N^--N}{T_N}&\le A_-\le\frac{S_N^-}{T_N},\label{rrm:eq:floorcert}\\
 \frac{S_N^+}{T_N}&\le A_+\le\frac{S_N^++N}{T_N}.\label{rrm:eq:ceilcert}
\end{align}
These certificates need only the integer trajectory and a positive rational Laguerre sum. They do not involve fitting an asymptotic curve, approximating $E_1$, iterating the decaying solution forward, or assuming that fractional parts are uniformly distributed.

For an integer-only normalization put
\begin{equation}\label{rrm:eq:Fdef}
 F_n=(n-1)!P_n.
\end{equation}
The recurrence becomes
\begin{equation}\label{rrm:eq:Frec}
 F_1=1,\quad F_2=2,\qquad
 F_{n+1}=2nF_n-(n-1)^2F_{n-1},
\end{equation}
and
\begin{equation}\label{rrm:eq:Tinteger}
 T_N=\frac{F_{N+1}-NF_N}{(N-1)!}.
\end{equation}
The main certificate uses \eqref{rrm:eq:Frec}. An independent replay uses the positive finite sums
\begin{align}
 F_N&=(N-1)!\sum_{j=0}^{N-1}\frac{\binom{N-1}{j}}{j!},\label{rrm:eq:positiveP}\\
 (N-1)!T_N&=N!\sum_{j=1}^{N}\frac{\binom{N-1}{j-1}}{j!}.\label{rrm:eq:positiveT}
\end{align}
The second identity follows by summing the finite formula for $P_n$ and using the hockey-stick binomial identity, or by expanding \eqref{rrm:eq:T}. Successive scaled terms are integers, so the replay can use exact integer divisions, explicitly checking every divisibility condition. This is a genuinely distinct computation of the normalization, rather than rerunning the same three-term recurrence under another name.

\section{Certified constants and numerical interpretation}\label{rrm:sec:numerics}
At $N=10000$, exact arithmetic in \eqref{rrm:eq:floorcert}--\eqref{rrm:eq:ceilcert} gives the following outward decimal intervals. Each displayed endpoint has 70 digits after the decimal point. To fit the page, each decimal is split after its first 35 fractional digits; the two strings on a line are to be concatenated.

\begin{center}\small
\setlength{\tabcolsep}{3pt}
\begin{tabular}{clll}
\toprule
Constant&End&First part&Continuation\\\midrule
$A_-$&lower&\texttt{0.73477794810835571221569276182082588}&\texttt{27018901106899598956984647699233851}\\
&upper&\texttt{0.73477794810835571221569276182082588}&\texttt{27018901106899598956984647699233852}\\[3pt]
$A_+$&lower&\texttt{1.29427700990317613682377668104327275}&\texttt{41541533342045100366386468876287149}\\
&upper&\texttt{1.29427700990317613682377668104327275}&\texttt{41541533342045100366386468876287150}\\\bottomrule
\end{tabular}
\end{center}

For readable full-precision records without line wrapping, the archive contains the exact generated certificate in JSON. The converted constants start as follows:
\begin{align}
 C_-&=0.445665353608456118285630970456186510059368576678\ldots,\notag\\
 c_-&=0.125719875127009200981669798844208976385113060072\ldots,\label{rrm:eq:decimalminus}\\
 C_+&=0.785018688667468005119788602907966565182706975887\ldots,\notag\\
 c_+&=0.221449683518252260781859024123926290928183228907\ldots.\label{rrm:eq:decimalplus}
\end{align}
These values agree with Kot\v e\v sovec's 2024 numerical observations in the two OEIS records. Agreement is a consistency check; the proof is supplied by the exact inequalities.

The common unrounded rational width $N/T_N$ is less than
\begin{equation}\label{rrm:eq:tinywidth}
 8.0967911913136482\times10^{-84}.
\end{equation}
Outward rounding to 70 places produces intervals of width $10^{-70}$. The displayed interval for $A_-$ certifies 69 common truncated fractional digits. The displayed $A_+$ endpoints cross a carry in the last two places and share only 68 fractional digits. A separate exact check of the narrower rational bounds verifies $\lfloor10^{69}L\rfloor=\lfloor10^{69}U\rfloor$ for each of the six constants $A_\pm,C_\pm,c_\pm$. Thus all six have 69 certified truncated fractional digits, as recorded in \texttt{certificates/common\_truncations.json}; that stronger statement uses the narrow bounds, not just the printed 70-place intervals.

\subsection{Why conversion to the Bessel constants is rigorous}
The factor $e^{-1/2}$ and $\sqrt\pi$ are enclosed with rational arithmetic. For $m\ge1$, put
\[
 E_m=\sum_{j=0}^{m-1}\frac1{2^jj!},\qquad
 R_m=\frac1{2^mm!}\frac{2(m+1)}{2(m+1)-1}.
\]
The omitted positive terms have successive ratios at most $1/[2(m+1)]$, so
\begin{equation}\label{rrm:eq:expbound}
 E_m<e^{1/2}<E_m+R_m.
\end{equation}
For $k>1$, the alternating expansion for $\arctan(1/k)$ places its value between the sum through index $m-1$ and that sum plus the first omitted term. Machin's exact identity
\[
 \pi=16\arctan(1/5)-4\arctan(1/239)
\]
therefore gives a rational lower and upper bound on $\pi$. To enclose a positive rational square root, let $q>0$, $B=10^d$, and
\[
 z=\left\lfloor\sqrt{\left\lfloor qB^2\right\rfloor}\right\rfloor.
\]
Then $z/B\le\sqrt q<(z+1)/B$. This uses only integer square root and integer division.

The implementation uses 160 exponential terms, 150 arctangent terms, and a square-root scale $10^{150}$. If $[A_L,A_U]$ is an amplitude enclosure, $[E_L,E_U]$ encloses $e^{1/2}$, and $[R_L,R_U]$ encloses $2\sqrt\pi$, then
\begin{equation}\label{rrm:eq:convert}
 \frac{A_L}{E_U}\le C\le\frac{A_U}{E_L},\qquad
 \frac{A_L}{E_UR_U}\le c\le\frac{A_U}{E_LR_L}.
\end{equation}
Every denominator is positive. Integer division rounds lower endpoints downward and upper endpoints upward. No floating-point result enters these intervals. The common-truncation check also uses these rational bounds before outward decimal rounding.

\subsection{Finite observations are not asymptotic onset certificates}
The package checks both complete displayed OEIS prefixes and all 1000 terms in each official b-file. It also checks finite Laguerre sums and exact Casoratians. Optional 120-digit diagnostics compare the six-correction forward approximation and the inverse with exact integer terms. Those floating results are evidence about selected indices, not interval proofs of an all-$n$ remainder or a smallest log-concavity onset.

\section{Further consequences for the rounded sequences}\label{rrm:sec:consequences}
\subsection{Beyond all orders proportionality}
Apply \eqref{rrm:eq:bothshadow} to both sequences and cancel $P_n$:
\begin{equation}\label{rrm:eq:proportion}
 a_n^+-\frac{A_+}{A_-}a_n^-=\OO(\sqrt n).
\end{equation}
Since $A_->0$, division by $a_n^-$ gives
\begin{equation}\label{rrm:eq:ratioorbit}
 \frac{a_n^+}{a_n^-}=\frac{A_+}{A_-}
 +\OO\left(n^{3/4}e^{-2\sqrt n}\right).
\end{equation}
Thus the normalized rounded sequences agree to every algebraic asymptotic order despite having different discontinuous rounding errors at many steps.

\subsection{Ratios and increments}
Taking enough terms of \eqref{rrm:eq:allorders}, shifting $n$ to $n+1$, and dividing produces an expansion to each fixed order. In particular,
\begin{align}
 \frac{a_{n+1}^\pm}{a_n^\pm}
 &=1+n^{-1/2}+\frac14n^{-1}+\OO(n^{-3/2}),\label{rrm:eq:ratios}\\
 a_{n+1}^\pm-a_n^\pm&\sim\frac{a_n^\pm}{\sqrt n}.\label{rrm:eq:increments}
\end{align}
For the stated ratio accuracy, using the forward theorem at a sufficiently deep fixed order controls the shifted remainders directly. One need not differentiate a big-$O$ expression.

\subsection{Eventual strict log concavity}
\begin{proposition}\label{rrm:prop:logconcavity}
Both rounded sequences are strictly log-concave for all sufficiently large $n$:
\begin{equation}\label{rrm:eq:logconcave}
 (a_n^\pm)^2>a_{n-1}^\pm a_{n+1}^\pm.
\end{equation}
No effective first such index is asserted.
\end{proposition}
\begin{proof}
Take the forward expansion through $c_3$, so the pointwise relative remainder is $\OO(n^{-2})$. Since its bracket tends to one, the logarithm has the form
\[
 \log a_n^\pm=\log c_\pm+2\sqrt n-\frac14\log n
 +q_1n^{-1/2}+q_2n^{-1}+q_3n^{-3/2}+\OO(n^{-2})
\]
for fixed rational $q_j$. The three shifted values of the last remainder still contribute only $\OO(n^{-2})$ to the central second difference. The explicitly written smooth terms can be differenced by Taylor's theorem; in particular the contribution of $2\sqrt n$ is $-1/(2n^{3/2})+\OO(n^{-7/2})$, and all other contributions are $\OO(n^{-2})$. Hence
\begin{equation}\label{rrm:eq:secondlog}
 \log a_{n+1}^\pm-2\log a_n^\pm+\log a_{n-1}^\pm
 =-\frac1{2n^{3/2}}+\OO(n^{-2})<0
\end{equation}
eventually. Exponentiating yields the claim.
\end{proof}

The proof uses a sufficiently small pointwise remainder before differencing. It does not assume that a nonsmooth rounding residual can itself be differentiated.

\section{All fixed orders of the range inverse}\label{rrm:sec:inverse}
Fix one of the two rounded sequences and suppress its sign. Write
\begin{equation}\label{rrm:eq:Dt}
 D=c\sqrt2=\frac{Ae^{-1/2}}{\sqrt{2\pi}},\qquad t=2\sqrt n.
\end{equation}
Then \eqref{rrm:eq:allorders} reads
\begin{equation}\label{rrm:eq:tforward}
 a_n=De^tt^{-1/2}
 \left\{1+\sum_{j=1}^M2^jc_jt^{-j}+\OO_M(t^{-M-1})\right\}.
\end{equation}
Taking the logarithm defines rational coefficients $b_j$ by
\begin{equation}\label{rrm:eq:blog}
 \sum_{j\ge1}b_jz^j
 =\log\left(1+\sum_{j\ge1}2^jc_jz^j\right).
\end{equation}
For every fixed $M$,
\begin{equation}\label{rrm:eq:logforward}
 \log(a_n/D)=t-\frac12\log t+
 \sum_{j=1}^M b_jt^{-j}+\OO_M(t^{-M-1}),
\end{equation}
with
\begin{equation}\label{rrm:eq:bvalues}
 b_1=-\frac{17}{24},\quad b_2=\frac5{16},\quad
 b_3=\frac{277}{1920},\quad b_4=\frac{21}{128},\quad
 b_5=\frac{85009}{107520}.
\end{equation}

\subsection{The large Lambert branch}
For sufficiently large $Y>0$, let $t_0$ be the large positive solution of $Y=De^{t_0}t_0^{-1/2}$. Squaring and rearranging gives
\[
 (-2t_0)e^{-2t_0}=-\frac{2D^2}{Y^2},
\]
so
\begin{equation}\label{rrm:eq:tzero}
 t_0=-\frac12W_{-1}\left(-\frac{2D^2}{Y^2}\right).
\end{equation}
Here $W_{-1}$ is the real branch tending to $-\infty$ as its argument tends to zero from below. The principal real branch $W_0$ would give a small root tending to zero and is not the inverse relevant to $n\to\infty$. This branch selection is part of the formula.

\subsection{Formal inversion and displayed coefficients}
Put $t=t_0+\delta$ and $v=1/t_0$. Subtract the defining leading equation for $t_0$ from \eqref{rrm:eq:logforward}. The formal equation is
\begin{equation}\label{rrm:eq:deltagen}
 \delta-\frac12\log(1+v\delta)
 +\sum_{j\ge1}b_jv^j(1+v\delta)^{-j}=0.
\end{equation}
Its derivative with respect to $\delta$ at $(v,\delta)=(0,0)$ equals one. Consequently there is a unique formal solution $\delta\in v\mathbb Q[[v]]$. Explicitly, if $d_1,\ldots,d_{r-1}$ have been found and $\delta_{<r}=\sum_{j<r}d_jv^j$, then $d_r$ is the negative coefficient of $v^r$ in the left side of \eqref{rrm:eq:deltagen} after substituting $\delta_{<r}$. The unknown $d_r$ appears only with coefficient one at that order.

The first terms are
\begin{align}
 \delta={}&\frac{17}{24t_0}+\frac1{24t_0^2}
 -\frac{3601}{5760t_0^3}-\frac{17}{90t_0^4}
 +\frac{223543}{967680t_0^5}+\OO(t_0^{-6}).
 \label{rrm:eq:delta}
\end{align}
Squaring $t/2$ gives the following inverse along the exact range.

\begin{theorem}[Range inverse]\label{rrm:thm:rangeinverse}
If $Y=a_n$ and $t_0$ is given by \eqref{rrm:eq:tzero} using the amplitude of that same sequence, then
\begin{align}
 n={}&\frac{t_0^2}{4}+\frac{17}{48}+\frac1{48t_0}
 -\frac{539}{2880t_0^2}-\frac{51}{640t_0^3}
 -\frac{25517}{241920t_0^4}+\OO(t_0^{-5}).
 \label{rrm:eq:rangeinverse}
\end{align}
The procedure \eqref{rrm:eq:blog}--\eqref{rrm:eq:deltagen} gives a justified expansion to every fixed inverse order. In particular any truncation whose proved error is $o(1)$ recovers the integer $n$ by nearest-integer rounding for all sufficiently large values in the exact sequence range.
\end{theorem}

\begin{proof}
We give a remainder argument, since formal reversion by itself is insufficient. Truncate \eqref{rrm:eq:logforward} at a fixed order $J$ and denote the smooth function on its right, without its remainder, by
\[
 f_J(t)=t-\tfrac12\log t+\sum_{j=1}^Jb_jt^{-j}.
\]
Then $f_J'(t)=1-1/(2t)+\OO_J(t^{-2})$, which is bounded below by $1/2$ for all sufficiently large $t$. The exact value $t=2\sqrt n$ satisfies
\[
 f_J(t)=\log(Y/D)+\OO_J(t^{-J-1}).
\]
First, the leading relation shows $t\sim t_0$. Substitution of the recursively constructed finite series $t_0+\delta_J$ into $f_J$ leaves a residual of order $\OO(t_0^{-J-1})$. The mean-value theorem therefore bounds $t-(t_0+\delta_J)$ by that same order. Squaring loses one power because $t+t_0+\delta_J=\OO(t_0)$, yielding an $n$-error of order $\OO(t_0^{-J})$.

With $J=5$, the coefficients \eqref{rrm:eq:bvalues} give \eqref{rrm:eq:delta}, and direct squaring gives \eqref{rrm:eq:rangeinverse}. Taking larger fixed $J$ proves every fixed inverse order. An error tending to zero is eventually less than $1/2$, which proves the nearest-integer statement on the exact range.
\end{proof}

The theorem uses the exact constant $D$. A decimal approximation held fixed while $Y\to\infty$ introduces an additional error; arbitrary-precision numerical evaluation must propagate the amplitude uncertainty if an executable inverse certificate is required. The all-order theorem does not identify an effective first index for nearest-integer recovery.

\section{Arbitrary thresholds and the staircase obstruction}\label{rrm:sec:threshold}
For arbitrary real $Y$, define
\begin{equation}\label{rrm:eq:threshold}
 N(Y)=\min\{n\ge1:a_n\ge Y\}.
\end{equation}
The sequence is strictly increasing and unbounded, so $N(Y)$ is well defined for every $Y$. Along its exact range, $N(a_n)=n$. Between successive terms it is a staircase, not a smooth inverse.

A continuous function cannot approximate this staircase with a universally valid additive $o(1)$ error over all sufficiently large real $Y$: at each large discontinuity its one-sided limits would have to approximate two integer values differing by one. Thus the range statement cannot simply be promoted to a single smooth threshold formula.

\subsection{Two monotone comparison functions}
For fixed $M\ge0$, let
\begin{equation}\label{rrm:eq:carrier}
 F_M(x)=ce^{2\sqrt x}x^{-1/4}\sum_{j=0}^Mc_jx^{-j/2}.
\end{equation}
It is positive and increasing for all sufficiently large real $x$. The forward theorem supplies constants $K>0$ and $n_0$ such that for integers $n\ge n_0$,
\begin{equation}\label{rrm:eq:carrierenvelope}
 F_M(n)(1-Kn^{-(M+1)/2})\le a_n
 \le F_M(n)(1+Kn^{-(M+1)/2}).
\end{equation}
The factors are positive and both comparison functions are increasing on a sufficiently far real half-line. Write
\[
 U_M(x)=F_M(x)(1+Kx^{-(M+1)/2}),\qquad
 L_M(x)=F_M(x)(1-Kx^{-(M+1)/2}).
\]
Let $x_+(Y)=U_M^{-1}(Y)$ and $x_-(Y)=L_M^{-1}(Y)$ on these increasing branches; the subscripts refer to the sign of the forward error. In particular $x_+(Y)\le x_-(Y)$.

\begin{theorem}[Two-ceiling threshold bounds]\label{rrm:thm:threshold}
For all sufficiently large $Y$,
\begin{equation}\label{rrm:eq:ceilbounds}
 \left\lceil x_+(Y)\right\rceil\le N(Y)
 \le\left\lceil x_-(Y)\right\rceil,
\end{equation}
and
\begin{equation}\label{rrm:eq:rootwidth}
 x_-(Y)-x_+(Y)=\OO(t_0^{-M}).
\end{equation}
For each fixed $M\ge1$ the two ceilings therefore eventually differ by at most one. For $M=0$ the root width is only $\OO(1)$, which does not imply adjacent candidate integers.
\end{theorem}
\begin{proof}
Take $Y$ large enough that both roots and all relevant indices lie beyond the monotonicity and forward-bound onset; also take $Y>\max_{n<n_0}a_n$. If an integer $n<\lceil x_+(Y)\rceil$, then $n<x_+(Y)$ and $a_n\le U_M(n)<Y$ for every relevant $n$. Thus $N(Y)\ge\lceil x_+(Y)\rceil$. At $n=\lceil x_-(Y)\rceil$, monotonicity gives $a_n\ge L_M(n)\ge Y$, proving the other side.

The two forward logarithms differ by $\OO(x^{-(M+1)/2})$. Their derivatives with respect to $x$ are $x^{-1/2}+\OO(x^{-1})$, uniformly on the relevant far half-line. Since the roots satisfy $x_\pm\sim t_0^2/4$, the mean-value theorem gives a root separation $\OO(t_0\cdot t_0^{-M-1})=\OO(t_0^{-M})$. For $M\ge1$ it is eventually less than one, so the two ceilings differ by at most one.
\end{proof}

\subsection{What the comparison functions do to inverse coefficients}
The logarithm of the additional factor $1\pm Kx^{-(M+1)/2}$ begins at order $t^{-(M+1)}$ after $t=2\sqrt x$. Its effect on the inverse $t$ begins at order $t_0^{-(M+1)}$, and its effect on $x=t^2/4$ begins at order $t_0^{-M}$. Therefore each root has its own formal inverse, obtained from \eqref{rrm:eq:deltagen} using the coefficients of that comparison function. The coefficients agree with the range inverse strictly before the first envelope-error order; the $\pm K$ terms can change coefficients at that order.

It is not valid to assign both roots all the same displayed range-inverse coefficients and append a smaller error. Nor does the theorem assert that taking one ceiling of a finite asymptotic truncation always gives the threshold. When actual forward constants $K,n_0$ are supplied and both computed ceilings coincide, \eqref{rrm:eq:ceilbounds} certifies the exact answer. Here their existence is proved but no numerical $K,n_0$ is certified, so the arbitrary-threshold statement is an asymptotic theorem rather than an executable finite threshold certificate.

\begin{remark}[{[write]} The inverses against the transseries volume, 7~October 2026]\label{rrm:rem:transseries}
The volume is
\path{Analysis/Transseries/docs/series-and-transseries/Transseries_And_Inversion/transseries_and_inversion.tex};
its symbols carry the subscript ``vol''. Statement by statement, for either
rounded sequence:
\begin{enumerate}[label=(\alph*)]
\item \emph{Instance.} $N(Y)$ of~\eqref{rrm:eq:threshold} is the integer
 staircase of \file{p0:def:three-inverses}(1) with $A_n=a_n$ and
 $n_{1,\mathrm{vol}}=1$ (the orbit is strictly increasing from $n=1$), so
 \file{p0:thm:staircase}(1) gives $N(Y)=\lceil\nu_{\mathcal A}(Y)\rceil$ for
 every admissible interpolation $\mathcal A_{\mathrm{vol}}$.
\item \emph{Analogue.} Theorem~\ref{rrm:thm:threshold} is a direct
 two-ceiling bracket at the integers about the roots of two comparison
 functions that do not interpolate the sequence: an analogue of
 \file{p0:thm:staircase}(2), not an instance.
\item \emph{Instance.} The nearest-integer recovery in
 Theorem~\ref{rrm:thm:rangeinverse} is \file{p0:thm:staircase}(3): its
 hypothesis involves only $\nu_{\mathcal A}(a_n)=n$, true for every admissible
 interpolation (the piecewise-linear one, say), and the truncated inverse
 approximates $n$ within $1/2$ for large $n$.
\item \emph{Instance, with the branch rule.} The Lambert root~\eqref{rrm:eq:tzero}
 solves $t_0-\tfrac12\log t_0=\log(Y/D)$, which is \file{p0:thm:lambert-core}
 with $a_{\mathrm{vol}}=1$, $b_{\mathrm{vol}}=-\tfrac12$ and
 $L_{\mathrm{vol}}=\log(Y/D)$: there $z(L_{\mathrm{vol}})=-2e^{-2\log(Y/D)}=-2D^2/Y^2$,
 and the branch rule for $b_{\mathrm{vol}}<0$ gives the large solution
 $X=(b_{\mathrm{vol}}/a_{\mathrm{vol}})W_{-1}(z)$, which is~\eqref{rrm:eq:tzero}
 with its $W_{-1}$ (for $\log(Y/D)>L_c=\tfrac12(1+\log2)$).
\item \emph{Formal instance.} Equation~\eqref{rrm:eq:deltagen} is
 \file{p0:thm:core-reversion} with $\Lambda_{\mathrm{vol}}=1$,
 $h_{\mathrm{vol}}=0$, $t_{\mathrm{vol}}=v$, $E_{\mathrm{vol}}=\delta$ and
 $f_{\mathrm{vol}}(t,u)=-\tfrac12\log(1+tu)+\sum_{j\ge1}b_jt^j(1+tu)^{-j}\in
 t\mathbb Q[[t,u]]$; its triangular recurrence is the source's. The analytic
 remainder argument in the proof of Theorem~\ref{rrm:thm:rangeinverse} is the
 source's own.
\end{enumerate}
No novelty is claimed for any of these inversions. The staircase arithmetic
is formalized generically (for an arbitrary strictly monotone function) as
\file{Fabius.staircase_ceil} in
\path{Analysis/FabiusFunction/Lean/FabiusFunction/StaircaseInversion.lean};
nothing about $a_n^\pm$ is formalized. (Corrected after the independent
check of 7~October 2026: \file{Fabius.staircase_ceil} is the arithmetic of
\file{p0:thm:staircase}(1), used in~(a). The nearest-integer step of~(c),
\file{p0:thm:staircase}(3), is \file{Fabius.staircase_round} in the same
file: $|x-n|\le\eta<1/2$ implies $\lfloor x+1/2\rfloor=n$.)
\end{remark}

\section{Reproducibility and the meaning of the checks}\label{rrm:sec:repro}
The companion archive contains this manuscript, the compiled report, documented Python source, the attributed official 1000-term b-files, and generated certificates. The mandatory computations use only Python's standard library. A TeX installation is needed only to rebuild the PDF.

\subsection{Mandatory exact computations}
\begin{enumerate}
\item \texttt{code/certify\_amplitudes.py} generates both integer trajectories to $N=10000$, computes the growing normalization using \eqref{rrm:eq:Frec}, proves the rational amplitude enclosures, encloses the exponential and square-root conversion factors, and rounds the resulting endpoints outward.
\item \texttt{code/check\_exact.py} replays both official 1000-term b-files and stored prefixes, checks the rational finite Laguerre sum against the recurrence through index 150, and checks exact Casoratians through index 150. For the latter, $Q_n$ is represented algebraically as $r_nQ_1+s_n$ with rational $r_n,s_n$; no transcendental approximation is necessary.
\item \texttt{code/audit\_positive\_sum.py} independently computes the normalization at $N=10000$ from \eqref{rrm:eq:positiveP}--\eqref{rrm:eq:positiveT} and replays the trajectories with integer division and remainder. It checks the amplitude brackets and their width against the generated certificate.
\item \texttt{code/formal\_series.py} generates rational forward and inverse coefficients with finite truncated-series algorithms and verifies the displayed coefficient identities. These algebraic checks identify coefficients; Perron's theorem supplies the analytic existence and remainder bounds.
\item \texttt{code/test\_mathematical\_guards.py} exercises meaningful corrupted inputs, so a broken prefix, interval, normalization, or coefficient cannot pass merely because a program runs to completion. Critical checks use explicit exceptions and remain active under optimized Python.
\end{enumerate}
The additional \texttt{code/check\_common\_truncations.py} verifies the 69-place truncations of all six constants from the narrow rational bounds and records the common-prefix lengths of the printed decimal endpoints separately. The reproduction driver runs mandatory checks in ordinary and optimized modes and compares their deterministic outputs with the frozen reference files. The manifest verifier checks the exact package inventory and file hashes. The build system refuses existing output destinations and isolates temporary TeX files and caches.

\subsection{Build and replay}
From an extracted archive, the principal commands are
\begin{quote}\ttfamily\small
python3 verify\_manifest.py\\
python3 reproduce.py\\
python3 build.py --output /absolute/path/to/new-output
\end{quote}
The last destination must be a new directory outside the extracted package, with an existing parent. The builder writes \texttt{Report186.pdf}, \texttt{Report186.tex}, and \texttt{Report186\_code.zip}, plus an artifact hash receipt. It uses fixed PDF metadata and ZIP timestamps, sorted inventories, disabled shell escape, and isolated temporary storage. Byte identity is promised only for the same source and installed Python/TeX stack; different TeX versions or fonts can change PDF bytes.

The package includes a detailed reproducibility guide with exact command options, provenance, test scope, optional dependencies, and the distinction between arithmetic certificates and diagnostics. Rebuilding does not require network access or alter any OEIS or other external record.

\subsection{Optional numerical and symbolic checks}
Optional high-precision diagnostics use mpmath at 120 decimal digits to inspect shadow residuals, six-correction relative errors, and range-inverse errors. Optional SymPy scripts rederive coefficients independently. These are supplementary experiments. They are not inputs to the exact amplitude proof, and numerical agreement is not promoted to an interval theorem.

\section{Attribution and bounded source review}\label{rrm:sec:sources}
The official defining records and b-files were inspected on 3 October 2026. The floor record was revision 19 dated 12 October 2024; the ceiling record was revision 20 dated 30 September 2026. Both records explicitly asked for a proof of the modified-Bessel equivalent at the time checked. The 2024 high-precision numerical constants are attributed in those records to V\'aclav Kot\v e\v sovec. Ulrich Schimke is the sequence author, and Harry J. Smith is credited for the linked 1000-term tables.
\begin{writenote}[Added 7 October 2026]
The write read the live records on 7~October 2026 (A065094 revision \#19,
A065095 revision \#20, unchanged): the question ``Can someone prove this?''
still stands in both, and Corollary~\ref{rrm:cor:oeis} answers it
(Remark~\ref{rrm:rem:oeis}).
\end{writenote}

The definitions are unambiguous despite two ancillary transcription inconsistencies: one floor-record comment writes its increment without the floor, and one ceiling-record example writes a floor symbol while evaluating the ceiling. The defining names, formulas, and complete verified integer prefixes agree with \eqref{rrm:eq:rounded}. No external correction or edit was made.

The unrounded sequence A376995 already identifies $L_{n-1}(-1)$ and its leading equivalent. We explicitly use that classical model. For the all-order input, the actual Dea\~no--Huertas--Marcell\'an paper was read at its normalization, Perron formula (3), negative-axis convention, Theorem 2, and formulas (24)--(27). Its all-order Laguerre expansion and coefficient methods are prior work. The first two coefficients were checked with its formula (27), including the degree shift. The standard modified-Bessel leading equivalent is also classical \cite{dlmf}.

A bounded search using the exact OEIS identifiers and rounded-mean, Laguerre-recurrence, and Bessel-growth terms did not reveal a prior proof of these particular rounded recurrences. Related models that were actually inspected concern exponentially weighted action inverses, balanced multiset words with fixed-degree Laguerre kernels, and digit-restricted partition products. Those are different models, although asymptotic inversion and Bessel comparison are shared established methods. This limited review is not an exhaustive literature search and does not establish global priority.

The self-contained content here is the positive minimal-solution construction, normalized Green identity, sharp bounded-forcing envelope, nested finite rational certificates, and their application to the two discontinuous recurrences. Classical results and observed constants are credited where used. The archive does not redistribute downloaded full third-party research papers.

\section{Further questions and limitations}\label{rrm:sec:questions}
The exact formulas isolate several questions that remain outside the present results.
\begin{enumerate}
\item \emph{Actual residual scale.} The optimal uniform envelope is of order $\sqrt n$, but it is not known here whether either rounded residual $a_n^\pm-A_\pm P_n$ has that order, a smaller order, or a limiting normalized profile. The forcing extremes used to prove sharpness need not be rounded trajectories.
\item \emph{Fractional parts.} No equidistribution, limiting average, independence, or other statistical law is proved for $\{S_n^\pm/n\}$. Such a law cannot be inferred from the amplitude formula, whose weights are summable and emphasize early errors.
\item \emph{Effective asymptotic onsets.} The exact amplitude intervals are executable certificates. The all-order remainder, eventual log-concavity, range-inverse nearest-integer recovery, and threshold comparison functions have no certified numerical onset here. Producing usable constants would require additional effective control of the classical remainder and its elementary transformations.
\item \emph{General rounding rules.} The bounded-forcing theorem already applies to nearest-integer rounding and any uniformly bounded deterministic correction. A separate positivity argument may be needed for the amplitude under a new initial condition or rule. The theorem permits $A=0$ in the general class, so positive exponential growth must not be assumed without that check.
\item \emph{Zero amplitude and constant forcing.} If every error equals $r$, the exact identity gives $a_n=(\alpha+r)P_n-r$. In particular $r=-\alpha$ gives a constant trajectory and zero amplitude. This example checks both the mass-one normalization and the necessity of positivity assumptions in relative asymptotics.
\item \emph{Dependence on errors.} For two forcings with the same initial value, $|A(\eps)-A(\eta)|\le\sup_k|\eps_k-\eta_k|$. If they agree through time $m-1$ and differ by at most $B$ afterward, the sharper bound is $|A(\eps)-A(\eta)|\le B(m-1)Q_{m-1}$ for $m\ge2$. These follow from the positive weights and exact tail; they quantify stability without a distributional model.
\end{enumerate}
\begin{writenote}[Added 7 October 2026, under Vladimir's standing rule of 4 October 2026]
The six questions above are the source's, credited to it and kept as stated.
The write found no unproved claim of the source outside them and no wrong
one. For Question~1 it computed a diagnostic, not a theorem: with the
certified 69-digit amplitudes, the actual residuals $r_n^\pm=a_n^\pm-A_\pm P_n$
satisfy $|r_n^-|\le3.30$ (largest at $n=946$) and $|r_n^+|\le4.08$ (largest
at $n=1528$) for $n\le2500$, while the envelope there is about
$\sqrt n/2\approx25$; their averages over $(n/2,n]$ lie between $0.43$ and
$0.69$ (floor) and between $-0.65$ and $-0.37$ (ceiling) for
$n=100,500,1000,1500,2000,2500$. No growth rate is inferred
from these values.
\end{writenote}

The central conclusion is that bounded rounding selects a scalar amplitude while preserving the complete classical algebraic asymptotic expansion. Its selection is both analytically explicit and certifiable from finite integer data.

\begin{writenote}[Independent check of the write, 7 October 2026]
An adversarial check made by the intake after the write (commit
\file{1ebd8c33c}) re-derived the statements marked [write] with its own code,
after fetching again A065094 (revision~\#19), A065095 (\#20), A376995
(\#16), the names of A160617 and A160618, and both b-files.
\emph{Remark~\ref{rrm:rem:oeis}.} Every quotation, revision, date and author
line confirmed (Kot\v{e}\v{s}ovec's comment is quoted in two fragments; the
elided words are ``It follows that the constant above is equal to''); both
b-files are byte-identical to the shipped ones and equal the check's orbits
in all $2000$ terms. \emph{The conjectures.} The proof chain of
Corollary~\ref{rrm:cor:oeis} re-read; by a route different from the
brackets~\eqref{rrm:eq:bracket}, namely $A=1+\sum_kw_k\eps_k$
of~\eqref{rrm:eq:A} with $w_k=k(Q_k-Q_{k+1})$ from a Miller backward
recurrence for the minimal solution normalized by $Q_1=eE_1(1)$, the check
gets $A_\pm$ to $80$ digits, agreeing with the bracket at $N=10000$ to
$3\cdot10^{-82}$; the four tabulated endpoints are the outward $70$-place
roundings of the exact brackets, the 69-digit common truncations hold, and
all four of Kot\v{e}\v{s}ovec's decimals ($50$, $48$, $49$, $47$ places) and
the source's four $48$-place values are truncations; $N/T_N<8.0967911913136482\cdot10^{-84}$.
\emph{Coefficients.} $c_1,\dots,c_6$ re-derived by an ansatz in the three-term
recurrence~\eqref{rrm:eq:recP} itself (not through the operator
of~\eqref{rrm:eq:operator}), $b_1,\dots,b_5$, the five coefficients
of~\eqref{rrm:eq:delta}, the range inverse~\eqref{rrm:eq:rangeinverse}, the
prefactor shift and $c_1=d_1-1$, $c_2=d_2-d_1+3/4$, and the ratio
expansion~\eqref{rrm:eq:ratios}: all agree. $T_N$ by the positive
sum~\eqref{rrm:eq:positiveT} equals~\eqref{rrm:eq:Tinteger} at $N=150$ and
$3000$, and the $N=3000$ brackets agree with those at $N=10000$ to $40$
places. \emph{The residual diagnostic} of Section~\ref{rrm:sec:questions}
reproduced ($3.294$ at $n=946$, $4.073$ at $n=1528$, $H_{2500}=24.50$, the
twelve averages), and both orbits satisfy the
envelopes~\eqref{rrm:eq:floorenv}--\eqref{rrm:eq:ceilenv} for $n\le2500$.
\emph{Remark~\ref{rrm:rem:transseries}.} (a)--(e) re-derived against the
volume: the instance claims hold, including $z(L_{\mathrm{vol}})=-2D^2/Y^2$
and $L_c=\tfrac12(1+\log2)$ for the branch rule; one correction, made by a
dated note there: the Lean declaration for~(c) is
\file{Fabius.staircase_round}, not \file{Fabius.staircase_ceil}.
\emph{Provenance.} Archive facts, the staged files (byte-identical), the 101
delivered label numbers and 63 references (aux of a build of the delivered
text, 19 pages) and Route~B (six certificates byte-identical, normal equal to
optimized, 20~s) confirmed. No error was found in the source's proofs.
\end{writenote}

\begin{thebibliography}{9}
\bibitem{oeisfloor}
The OEIS Foundation Inc., \emph{A065094}, sequence contributed by Ulrich Schimke. Floor running-mean recurrence; numerical asymptotic constants credited to V\'aclav Kot\v e\v sovec, 12 October 2024. Definition and record checked 3 October 2026 (and by the write on 7~October 2026, revision \#19).
\url{https://oeis.org/A065094}.

\bibitem{oeisceil}
The OEIS Foundation Inc., \emph{A065095}, sequence contributed by Ulrich Schimke. Ceiling running-mean recurrence; numerical asymptotic constants credited to V\'aclav Kot\v e\v sovec, 12 October 2024. Definition and record checked 3 October 2026 (and by the write on 7~October 2026, revision \#20).
\url{https://oeis.org/A065095}.

\bibitem{bfiles}
Harry J. Smith, tables of the first 1000 terms of A065094 and A065095, linked from the official OEIS entries.
\url{https://oeis.org/A065094/b065094.txt} and
\url{https://oeis.org/A065095/b065095.txt}.

\bibitem{oeisunrounded}
The OEIS Foundation Inc., \emph{A376995}, the unrounded Laguerre running-mean sequence. Checked 3 October 2026.
\url{https://oeis.org/A376995}.

\bibitem{dhm}
Alfredo Dea\~no, Edmundo J. Huertas, and Francisco Marcell\'an,
\emph{Strong and ratio asymptotics for Laguerre polynomials revisited},
arXiv:1301.4266v2. The classical Perron theorem is formula (3); the coefficient and branch checks used here are in Theorem 2 and formulas (24)--(27).
\url{https://arxiv.org/abs/1301.4266}.

\bibitem{dlmf}
F. W. J. Olver and others, editors, \emph{NIST Digital Library of Mathematical Functions}, Section 10.40, formula 10.40.1, modified-Bessel asymptotic expansion for fixed order and large positive argument.
\url{https://dlmf.nist.gov/10.40.E1}.
\end{thebibliography}
\end{document}
